%! Functions and Function Notation — Unit 1, Lesson 1.1 — AP Practice
% EXTRA CREDIT — optional, exactly TWO pages, placed between the notes and the graded homework.
% Shaped like the AP Precalculus exam: page 1 is Section I, five multiple-choice items with four
% options (A)–(D), each in context or on a pre-drawn display; page 2 is Section II, one
% multi-part free-response question with interpret-in-context parts. Its contexts (a reservoir,
% a shipping tariff) are used nowhere else in the lesson.
% Distractors are the lesson's errors on purpose: MC 1 reads f(x) as multiplication, MC 2 drops
% the parentheses on a negative input, MC 3 uses the rule that REACHES the boundary instead of
% the one that OWNS it (the lesson's crux), MC 4 swaps domain and range, MC 5 misses the
% coefficient when clearing the denominator.
%
% PAGE LOCKSTEP with the other half of the pair: the key marks the correct option in its
% LABEL (no text added to the line) and prose answers are \writespace{H}{…}, fixed height both
% ways. The explicit \newpage fixes the page break in both files.
\documentclass[10pt]{article}
\usepackage{precalculus-article}
\usepackage{precalculus-boxes}
\newcommand{\TallMath}[1]{$\displaystyle #1\rule[-1.4em]{0pt}{3.2em}$}

\begin{document}

\pageheader{Unit 1, Lesson 1.1}{AP Practice}

\vspace{0.12in}

\boxguard[8]
\begin{remindbox}
\small
\textbf{Extra credit --- optional.} These two pages are written in the style of the AP
Precalculus exam. Your graded homework is the last section of this packet; do it first. Every
answer choice below is a mistake someone really makes, so read all four before you choose.
\end{remindbox}

\vspace{0.12in}

\begin{headlinebox}{goldbg}
\small
\textbf{Section I --- Multiple Choice.} Circle the letter of the single best answer. No calculator.
\end{headlinebox}

\vspace{0.06in}

\begin{enumerate}[leftmargin=1.9em, itemsep=12pt]

\item The function $W$ gives the water level of a reservoir, in feet, $t$ days after June 1.
      Which of the following is the best interpretation of $W(10) = 42$?
      \begin{enumerate}[leftmargin=2.6em, itemsep=2pt]
        \item[(A)] After 42 days, the water level is 10 feet.
        \item[(B)] Ten days after June 1, the water level is 42 feet.
        \item[(C)] The water level rose 42 feet over the first 10 days.
        \item[(D)] $W$ multiplied by 10 is 42, so $W = 4.2$.
      \end{enumerate}

\item If $f(x) = x^2 - 4x$, what is the value of $f(-3)$?

      \smallskip
      \textbf{(A)} $-21$ \hspace{2.2em} \textbf{(B)} $-3$ \hspace{2.2em} \textbf{(C)} $3$ \hspace{2.2em} \textbf{(D)} $21$

\item The function $p$ is defined by $p(x) = 3x + 1$ for $x < 4$, and $p(x) = x^2 - 6$ for
      $x \ge 4$. What is the value of $p(4)$?

      \smallskip
      \textbf{(A)} $10$ \hspace{2.2em} \textbf{(B)} $13$ \hspace{2.2em} \textbf{(C)} $23$ \hspace{2.2em} \textbf{(D)} Both 10 and 13

\item \begin{minipage}[t]{0.56\linewidth}
      The graph of the function $h$ is shown, including both endpoints. Which of the following
      gives the domain and the range of $h$?
      \begin{enumerate}[leftmargin=2.6em, itemsep=3pt]
        \item[(A)] Domain $-4 \le x \le 3$; range $-2 \le y \le 4$
        \item[(B)] Domain $-2 \le x \le 4$; range $-4 \le y \le 3$
        \item[(C)] Domain $-4 \le x \le 3$; range $1 \le y \le 4$
        \item[(D)] Domain $-4 \le x \le 4$; range $-2 \le y \le 3$
      \end{enumerate}
      \end{minipage}\hfill
      \begin{minipage}[t]{0.38\linewidth}
      \vspace{-4pt}
      \centering
      \begin{tikzpicture}[x=0.42cm, y=0.36cm, baseline=(current bounding box.north)]
        \draw[linegray, very thin, step=1] (-5,-3) grid (4,5);
        \draw[->, thick] (-5,0) -- (4.5,0) node[right] {\scriptsize $x$};
        \draw[->, thick] (0,-3) -- (0,5.5) node[above] {\scriptsize $h(x)$};
        \foreach \x in {-4,-2,2} \node[below, font=\tiny] at (\x,0) {\x};
        \foreach \y in {-2,2,4} \node[left, font=\tiny] at (0,\y) {\y};
        \draw[very thick, plum] (-4,1) -- (-1,-2) -- (3,4);
        \fill[plum] (-4,1) circle (2.2pt);
        \fill[plum] (3,4) circle (2.2pt);
      \end{tikzpicture}
      \end{minipage}

\item What is the domain of $g(x) = \dfrac{5}{2x - 8}$?
      \begin{enumerate}[leftmargin=2.6em, itemsep=2pt]
        \item[(A)] All real numbers
        \item[(B)] All real numbers except $-4$
        \item[(C)] All real numbers except $4$
        \item[(D)] All real numbers except $8$
      \end{enumerate}

\end{enumerate}

\newpage

\begin{headlinebox}{goldbg}
\small
\textbf{Section II --- Free Response.} Show your reasoning. Answers about the shipping cost
must be written \emph{in context}, with units --- that is what earns credit on the AP exam.
\end{headlinebox}

\vspace{0.08in}

\begin{enumerate}[leftmargin=1.9em, itemsep=10pt, start=6]

\item A shipping company charges a flat rate for light packages and a per-pound rate above
      that. For a package weighing $w$ pounds, the cost in dollars is
      \[ C(w) = \begin{cases} 6 & 0 < w \le 2 \\[2pt] 6 + 2(w - 2) & 2 < w \le 10 \end{cases} \]

      \textbf{(a)} Find $C(2)$ and $C(5)$. Interpret $C(5)$ in the context of the problem, with
      units.
      \writespace{2.2cm}{}

      \textbf{(b)} A customer claims that because $2$ appears in both conditions, $C(2)$ has two
      possible values. Explain why the customer is wrong.
      \writespace{2.4cm}{}

      \textbf{(c)} Explain what would go wrong with the definition of $C$ if the first rule were
      written $0 < w < 2$ and the second were left as $2 < w \le 10$.
      \writespace{2.4cm}{}

      \textbf{(d)} State the domain of $C$, as an inequality and in bracket notation. Then
      explain whether $C(12)$ has a meaning in this context.
      \writespace{2.4cm}{}

      \textbf{(e)} The company adds a flat \$25 charge for any package over 10 pounds, up to a
      50-pound limit. Write the third rule, with its condition, as it would appear in $C$.
      \writespace{2.2cm}{}

\end{enumerate}

\end{document}
